\section*{Bab \ref{ch:b1:hotspots} --- Hotspot dan vulkanisme intralempeng}

\begin{solution}{exo:b1:hotspots:1}
(a) $4205 + 5274 = \qty{9479}{m}$, hampir \qty{9.5}{km}. Gunung api itu lebih berat daripada air yang digantikannya: lempeng membengkok ke bawah di bawah bebannya, sehingga dasarnya terletak lebih dalam daripada dasar jauh di sana, dan tinggi sebenarnya lebih besar. (b) $74000/0.1 = \num{740000}$ tahun dengan pasokan Kilauea; $74000/0.025 = \num{3.0e6}$ tahun dengan laju rata-rata rantai. (c) $0.12\times19 = \qty{2.3}{km^3}$, sekitar $\num{3e-5}$ dari Mauna Loa: tiga perseribu persen.
\end{solution}

\begin{solution}{exo:b1:hotspots:2}
(a) $0.92\times\pi/180 = \qty{0.0161}{rad.Ma^{-1}}$. (b) Dalam radian per juta tahun dikali kilometer, $\omega R\sin\Delta$ berada dalam kilometer per juta tahun, yaitu milimeter setahun: $0.016057\times6371 = 102.3$, sehingga $v = 102.3\sin\Delta$, yang memberi \qty{51.1}{mm.yr^{-1}} pada \ang{30}, \qty{88.6}{mm.yr^{-1}} pada \ang{60}, \qty{102.3}{mm.yr^{-1}} pada \ang{90} dan \qty{97.3}{mm.yr^{-1}} pada \ang{108.0}. (c) Lajunya sebanding dengan jarak $R\sin\Delta$ dari sumbu, terbesar pada \omterm{def:b1:spherical-astronomy:spherical-triangle}{lingkaran besar} \ang{90} dari kutub; melewatinya titik itu kembali mendekati sumbu, di sisi kutub yang berlawanan.
\end{solution}

\begin{solution}{exo:b1:hotspots:3}
(a) Deccan: $\num{e6}/\num{5e5} = \qty{2}{km}$. Atlantik Tengah: $\num{2e6}/\num{7e6} = \qty{0.29}{km}$, sekitar \qty{290}{m}. (b) Rata-rata untuk Deccan dekat dengan penampangnya, \qty{2.5}{km}: provinsi itu adalah dataran tebal yang menerus. (c) Luas Atlantik Tengah adalah luas seluruh wilayah, di atas empat benua, tempat sisa-sisanya ditemukan: lava, dike dan sill yang kini terpisah oleh samudra dan banyak terkikis. Batuan itu terkonsentrasi di beberapa tempat dan tidak ada di sebagian besar wilayah, sehingga rata-rata tidak mengatakan apa pun tentang satu tumpukan saja.
\end{solution}

\begin{solution}{exo:b1:hotspots:4}
(a) $\rho_m\alpha\,\Delta T = 3300\times\num{3.2e-5}\times250 =
\qty{26.4}{kg.m^{-3}}$. (b) Kolom tonjolan memiliki massa tambahan $(\rho_m - \rho_w)h = 2300\times690 = \qty{1.587e6}{kg.m^{-2}}$, yang harus diimbangi oleh lapisan ringan: $H = \num{1.587e6}/26.4 = \qty{6.0e4}{m}$, sekitar \qty{60}{km} batuan panas.
\end{solution}

\begin{solution}{exo:b1:hotspots:5}
(a) Di A ($\lambda = 0$, $\varphi = 0$), $\vec e = (0, 1, 0)$ dan
$\vec n = (0, 0, 1)$; dengan $\sin\ang{270} = -1$ dan $\cos\ang{270} = 0$,
$\vec t_A = -\vec e = (0, -1, 0)$. Di B ($\lambda = \ang{90}$),
$\vec e = (-1, 0, 0)$, $\vec n = (0, 0, 1)$ dan $\vec t_B = (1, 0, 0)$.
(b) $\vec t_A\times\vec t_B = (0, 0, 1)$. (c) Dengan $\vec p' = (0, 0, 1)$ dan
$\vec r_A = (1, 0, 0)$, $\vec p'\times\vec r_A = (0, 1, 0)$, dan
$(\vec p'\times\vec r_A)\cdot\vec t_A = -1 < 0$: kutubnya adalah
$-\vec p' = (0, 0, -1)$, Kutub Selatan. Gerak ke barat sepanjang khatulistiwa
searah jarum jam dilihat dari atas Kutub Utara, sehingga berlawanan arah jarum jam dilihat dari atas
Kutub Selatan. (d) Pada \ang{90} dari kutub, $\omega = 50/6371 =
\qty{7.85e-3}{rad.Ma^{-1}} = \qty{0.45}{\degree\per\Ma}$.
\end{solution}

\begin{solution}{exo:b1:hotspots:6}
(a) $\sum x_i v_i = 90000 + 240000 + 342000 = \num{672000}$ dan
$\sum x_i^2 = \num{7.0e7}$, sehingga $\omega = \qty{0.0096}{rad.Ma^{-1}}$, atau
\qty{0.55}{\degree\per\Ma}. (b) Prediksinya adalah 28.8, 48.0 dan
\qty{57.6}{mm.yr^{-1}}; residunya $+1.2$, $0.0$ dan $-0.6$
\unit{mm.yr^{-1}}. (c) Untuk rotasi kaku, laju sebanding dengan
jarak dari sumbu, dan titik pada sumbu tidak bergerak: tidak ada
perpotongan yang perlu dicocokkan.
\end{solution}

\begin{solution}{exo:b1:hotspots:7}
(a) $v = \qty{50e-3}{m}/\qty{3.156e7}{s} = \qty{1.58e-9}{m.s^{-1}}$, sehingga
$B = 2300\times\num{3e8}\times\num{1.58e-9} = \qty{1090}{kg.s^{-1}}$ dan
$Q = (910/\num{3.2e-5})\times1090 = \qty{3.1e10}{W}$. (b) Sepertiga dari
fluks Hawaii. (c) $B$ tumbuh dengan $A\,v$: plume di bawah lempeng yang lebih cepat
harus membuat ulang tonjolan yang sama dua kali lebih cepat, dan fluksnya dua kali lebih besar.
\end{solution}

\begin{solution}{exo:b1:hotspots:8}
(a) $350/(16.1 - 10.3) = 350/5.8 = \qty{60}{km.Ma^{-1}}$, yaitu
\qty{60}{mm.yr^{-1}}. (b) Jejak muda, \qty{29 \pm 5}{mm.yr^{-1}}, sesuai
dengan \qty{26 \pm 5}{mm.yr^{-1}} dari lempeng dalam ketidakpastiannya; bagian
tua bergerak lebih dari dua kali lebih cepat, $60 - 26 = 34$, hampir tujuh simpangan baku
ketidakpastian laju lempeng. (c) Ya: U\,\ang{54}\,T adalah \omterm{def:b1:spherical-astronomy:horizontal}{azimut}
\ang{54}, berlawanan dengan \ang{235} dalam satu derajat. Vulkanisme bergerak di atas
lempeng ke arah yang berlawanan dengan lempeng di atas \omterm{def:b1:hotspots:hotspot}{hotspot}. (d) Peregangan memindahkan
kaldera-kaldera yang lebih tua menjauh, memanjangkan jarak di antaranya dan karenanya
laju semu; lempeng itu sendiri tidak bergerak lebih cepat.
\end{solution}

\begin{solution}{exo:b1:hotspots:9}
(a) $\num{9.2e10}/\num{e12} = \qty{0.092}{W.m^{-2}}$, \qty{92}{mW.m^{-2}}.
(b) Sembilan hingga delapan belas kali kelebihan yang teramati, yang bahkan mungkin nol.
(c) $L^2/\kappa =
\num{e10}/\num{e-6} = \qty{e16}{s}$, sekitar \qty{320}{Ma}; lempeng menyeberangi
\qty{1000}{km} dalam $1000/92 = \qty{11}{Ma}$. Panas dari bawah sempat
menyeberangi hanya sebagian kecil lempeng selama ia berada di atas tonjolan: sebagian besar
panas plume tetap di dalam \omterm{def:b1:layered-earth:layers}{mantel} di bawah lempeng dan terbawa pergi bersamanya, dan permukaan hanya melihat sedikit darinya.
\end{solution}

\begin{solution}{exo:b1:hotspots:10}
(a) Sepuluh kali lelehan, \qty{e7}{km^3}, dengan jari-jari
$(3V/4\pi)^{1/3} = \qty{134}{km}$. (b) $\num{e7}/100 = \qty{e5}{km^2}$, sebuah piringan
$2\sqrt{\num{e5}/\pi} = \qty{357}{km}$ diameternya: delapan belas kali lebih kecil luasnya
daripada Trap. Kepalanya lebih besar atau meleleh lebih sedikit, atau ia menyebar lebih lebar dan
lebih tipis, atau lavanya mengalir jauh melampauinya. (c) \qty{1}{km^3} setahun, sepuluh
kali pasokan Kilauea saat ini. (d) Sebuah plume baru harus mengumpulkan cukup
batuan panas dan ringan untuk mendorong jalannya naik melalui \omterm{def:b1:layered-earth:layers}{mantel} kental, dan volume besar
naik jauh lebih cepat daripada volume kecil; ekor yang mengikutinya mengalir ke atas melalui
jalur panas dan lemah yang dibuka kepala, dan saluran sempit sudah cukup.
\end{solution}

\begin{solution}{exo:b1:hotspots:11}
(a) $\qty{40}{km.Ma^{-1}}\times(81 - 47)\,\unit{Ma} = \qty{1360}{km}$.
(b) Lempeng bergerak dengan $(-80\sin\ang{20}, 80\cos\ang{20}) = (-27.4, 75.2)$
dan \omterm{def:b1:hotspots:hotspot}{hotspot} pada titik $(0, -40)$: relatif terhadap \omterm{def:b1:hotspots:hotspot}{hotspot} lempeng bergerak dengan vektor
$(-27.4, 115.2)$ \unit{mm.yr^{-1}}, laju \qty{118.4}{mm.yr^{-1}} ke arah
$\ang{360} - \arctan(27.4/115.2) = \ang{346.6}$. (c) Jejak akan memberi
laju \qty{48}{\percent} terlalu tinggi dan arah \ang{6.6} terlalu jauh ke timur:
\omterm{def:b1:hotspots:hotspot}{hotspot} yang hanyut mencatat gerak lempeng relatif terhadap dirinya sendiri, bukan terhadap
\omterm{def:b1:layered-earth:layers}{mantel} dalam.
\end{solution}

\begin{solution}{exo:b1:hotspots:12}
(a) $\qty{3e6}{J.m^{-3}.K^{-1}}\times\qty{1300}{K}\times\qty{e15}{m^3} =
\qty{3.9e24}{J}$. (b) Dalam $\qty{3.16e13}{s}$, \qty{1.2e11}{W}: kira-kira sama dengan
panas yang dibawa plume Hawaii saat ini (1.3 kali), dan \qty{0.26}{\percent}
dari \qty{47}{TW} Bumi. (c) Batas bawah: kepala juga membawa naik panas berlebih dari
sembilan persepuluh batuannya yang tidak meleleh, dan memanaskan dasar \omterm{def:b1:layered-earth:lithosphere}{litosfer}, dan panas laten dikeluarkan.
\end{solution}

\begin{solution}{pb:b1:hotspots:1}
\textbf{1.} Di Hawaii ($\varphi = \ang{19.0}$, $\lambda = \ang{-155.2}$),
dengan $\sin\ang{304} = -0.8290$ dan $\cos\ang{304} = 0.5592$:
\[
\begin{aligned}
  \vec r_H &= (-0.8583,\ -0.3966,\ 0.3256), &
  \vec e &= (0.4195,\ -0.9078,\ 0),\\
  \vec n &= (0.2955,\ 0.1366,\ 0.9455), &
  \vec t_H &= (-0.1825,\ 0.8289,\ 0.5287).
\end{aligned}
\]

\textbf{2.} Di Louisville ($\varphi = \ang{-53.6}$, $\lambda = \ang{-140.6}$),
dengan $\sin\ang{316} = -0.6947$ dan $\cos\ang{316} = 0.7193$:
\[
\begin{aligned}
  \vec r_L &= (-0.4586,\ -0.3767,\ -0.8049), &
  \vec e &= (0.6347,\ -0.7727,\ 0),\\
  \vec n &= (-0.6220,\ -0.5109,\ 0.5934), &
  \vec t_L &= (-0.8883,\ 0.1693,\ 0.4269).
\end{aligned}
\]

\textbf{3.} $\lVert\vec t_H\rVert^2 = 0.0333 + 0.6871 + 0.2795 = 0.9999$ dan
$\vec t_H\cdot\vec r_H = 0.1566 - 0.3287 + 0.1721 = 0.0000$, sampai pembulatan.

\textbf{4.} $\vec t_H\times\vec t_L = (0.2643, -0.3917, 0.7054)$, dengan panjang
$0.8491$. Sudut antara jejak-jejak itu adalah $\arcsin0.8491 = \ang{58.1}$ (hasil kali
titiknya, $0.528$, positif, sehingga sudutnya lancip).

\textbf{5.} $\vec p' = (0.3113, -0.4613, 0.8308)$. Maka
$(\vec p'\times\vec r_H)\cdot\vec t_H = -0.98 < 0$: kutubnya adalah
$\vec p = (-0.3113, 0.4613, -0.8308)$.

\textbf{6.} $\arcsin(-0.8308) = \ang{-56.18}$ dan
$\operatorname{atan2}(0.4613, -0.3113) = \ang{124.01}$: kutub terletak pada
\ang{56.2}\,LS, \ang{124.0}\,BT, di Samudra Selatan selatan Australia.

\textbf{7.} $\vec p\cdot\vec r_H = 0.2672 - 0.1830 - 0.2705 = -0.1863$, sehingga
$\Delta_H = \ang{100.7}$; $\vec p\cdot\vec r_L = 0.1428 - 0.1738 + 0.6687 =
0.6377$, sehingga $\Delta_L = \ang{50.4}$.

\textbf{8.} Vektor $\vec p\times\vec r_L$ tegak lurus terhadap $\vec p$ dan
terhadap $\vec r_L$, seperti $\vec t_L$ (menurut konstruksi $\vec p$, dan karena jejak itu
horizontal); keduanya terletak pada garis yang sama, dan panjangnya adalah
$\sin\Delta_L$. Jadi $(\vec p\times\vec r_L)\cdot\vec t_L = \pm\sin\Delta_L$, dan
perhitungan memberi $+0.770 = \sin\ang{50.4}$: lempeng bergerak di Louisville
sepanjang $+\vec t_L$, ke arah jejak itu berjalan. Aturan tanda, jika diterapkan di Hawaii,
bisa memberi arah yang salah di Louisville; ternyata tidak, sehingga kedua arah direproduksi
tepat, \ang{304} dan \ang{316}.

\textbf{9.} $\omega = 67/(6371\sin\ang{50.38}) = 67/4908 =
\qty{0.01365}{rad.Ma^{-1}}$, yaitu \qty{0.782}{\degree\per\Ma}.

\textbf{10.} $360/0.782 = \qty{460}{Ma}$.

\textbf{11.} Aturan kosinus antara \ang{56.18}\,LS, \ang{124.01}\,BT dan
\ang{63.78}\,LS, \ang{112.48}\,BT memberi \ang{9.5}. Laju dua jejak,
\qty{0.78}{\degree\per\Ma}, \qty{15}{\percent} di bawah
\qty{0.92}{\degree\per\Ma} yang dicocokkan pada tiga belas laju.

\textbf{12.} Society \ang{73.6}, Pitcairn \ang{77.8}, Cobb \ang{134.7}.

\textbf{13.} $\omega R = 0.01365\times6371 = \qty{86.98}{mm.yr^{-1}}$, dikali
$\sin\Delta$: Society \qty{83.5}{mm.yr^{-1}}, Pitcairn
\qty{85.0}{mm.yr^{-1}}, Cobb \qty{61.8}{mm.yr^{-1}}.

\textbf{14.} Pitcairn sesuai: \qty{5}{mm.yr^{-1}} lebih kecil daripada yang terukur, dalam
ketidakpastiannya sebesar \qty{15}{mm.yr^{-1}}. Society \qty{25.5}{mm.yr^{-1}} terlalu lambat, 2.6
simpangan baku, dan Cobb \qty{18.8}{mm.yr^{-1}} terlalu cepat, lebih dari enam:
tidak satu pun yang sesuai.

\textbf{15.} Kutub tegak lurus terhadap kedua vektor jejak. Jika keduanya hampir sejajar,
hasil kali silangnya pendek, dan kesalahan beberapa derajat ke arah mana pun memutarnya
melalui sudut besar: kutubnya buruk ditentukan.

\textbf{16.} $v = \qty{92e-3}{m}/\qty{3.156e7}{s} = \qty{2.92e-9}{m.s^{-1}}$
dan $B = 2300\times\num{4.83e8}\times\num{2.92e-9} = \qty{3240}{kg.s^{-1}}$.

\textbf{17.} $Q = (910/\num{3.2e-5})\times3240 = \qty{9.2e10}{W}$,
\qty{0.20}{\percent} dari \qty{47}{TW}.

\textbf{18.} Punggungan membuat $21/67000 = \qty{3.1e-4}{km^3}$ setahun per kilometer;
$0.1/\num{3.1e-4} = \qty{320}{km}$ punggungan.

\textbf{19.} $74000/0.1 = \num{740000}$ tahun.

\textbf{20.} $\num{2e6}/\num{81e6} = \qty{0.025}{km^3}$ setahun: Kilauea saja hari ini
dipasok dengan empat kali keluaran rata-rata seluruh \omterm{def:b1:hotspots:hotspot}{hotspot}, yang hampir berada pada
puncaknya.

\textbf{21.} $v_H = 86.98\sin\ang{100.73} = 86.98\times0.9825 =
\qty{85.5}{mm.yr^{-1}}$.

\textbf{22.} $92 - 85.5 = \qty{6.5}{mm.yr^{-1}}$, 2.2 simpangan baku dari laju yang terukur:
kedua \omterm{def:b1:hotspots:hotspot}{hotspot} tidak sepenuhnya sepakat tentang gerak satu \omterm{def:b1:plate-kinematics:rigid-plate}{lempeng kaku} di atas kerangka tetap.

\textbf{23.} Umur memberi \qty{9.1}{cm} dan \qty{8.6}{cm} setahun, dirata-ratakan
atas seluruh rantai; \qty{9.2}{cm} setahun dari Morgan dan Morgan dekat dengan yang pertama,
dan prediksi Louisville, \qty{8.5}{cm} setahun, dekat dengan yang kedua. Keempatnya berada dalam
\qty{0.7}{cm} setahun satu sama lain: jejak, umur dan arah menceritakan kisah yang sama hingga
sekitar \qty{8}{\percent}.

\textbf{24.} Hingga satu desimal, jejak Louisville memberikan kepada Lempeng Pasifik di atas
\omterm{def:b1:hotspots:hotspot}{hotspot} Hawaii laju
\textbf{\qty[text-series-to-math=true]{8.5}{cm.yr^{-1}}}.
\end{solution}
